% Основной текст новой версии. Подключается из stage1_report.tex.
\section{Постановка задачи и статус согласования}
\subsection{Тема и цель работы}
Тема курсовой работы: «Информационная система поиска салонов красоты и онлайн-записи на услуги».

Пользователь, которому нужна стрижка или маникюр, может не знать ближайшие салоны, работающих там специалистов и цены. Платформа объединяет несколько салонов и позволяет найти их на карте или в списке, сравнить услуги, изучить фотографии и отзывы, выбрать мастера и свободное время, оформить либо отменить запись.

Цель системы --- обеспечить путь от поиска подходящего салона до записи и оценки состоявшегося посещения. Система предназначена для клиентов, мастеров, представителей салонов и администратора платформы. Термин «мастер» включает парикмахеров, специалистов по маникюру и других специалистов красоты.

Первый этап включает анализ предметной области, требования, модели и описания прецедентов и архитектуру. Прикладной код и физическая схема базы данных на данном этапе не разрабатываются.

\subsection{Предложения для согласования}
Тема, прецеденты и технологии являются предложением; факт утверждения преподавателем не предполагается. Требуется согласовать:
\begin{enumerate}
\item платформу нескольких салонов с четырьмя ролями и границы учебной версии;
\item карту и список, фильтры, страницы салонов и мастеров, фотографии, запись, отмену и отзывы;
\item Vue, Spring Boot со Spring MVC и PostgreSQL;
\item среду \texttt{helios}, защищённый доступ из браузера, картографическую подложку и место для изображений.
\end{enumerate}
Дата и результат согласования: \textit{заполняются после обсуждения с преподавателем}.

\section{Описание предметной области}
\subsection{Участники и процесс обслуживания}
Клиент ищет услугу в удобном месте и в пределах бюджета. Для выбора нужны адрес, цена, длительность, мастера, свободное время, примеры работ и отзывы. Разрозненный поиск и ручное согласование по телефону затрудняют сравнение и создают риск ошибки в расписании.

Представитель салона, далее также «владелец», регистрирует организацию, ведёт её описание, услуги, цены, состав мастеров, рабочие интервалы и фотографии. Он управляет только своим салоном. Платформа предоставляет ему публичную страницу и журнал записей.

Мастер имеет учётную запись и профессиональный профиль. Он просматривает своё расписание, размещает работы, отменяет при необходимости будущие записи к себе и отмечает результат посещения. Цены и рабочие интервалы в базовой версии устанавливает владелец.

Администратор платформы поддерживает общие справочники, проверяет заявки организаций, управляет доступом и модерирует отзывы и изображения. Он отделён от владельца и не ведёт повседневное расписание каждого салона.

\subsection{Карта, список и фильтрация}
Без фильтров список содержит все опубликованные активные салоны, зарегистрированные на платформе, с постраничным просмотром. Неопубликованные заявки и заблокированные организации в публичную выдачу не входят.

Пользователь выбирает район из справочника, область карты либо режим «Рядом со мной». Последний запрашивает разрешение браузера на определение положения и использует радиус, например 1, 3 или 5 километров. Географические режимы альтернативны; текущий режим явно показан. При отказе от геолокации остаются ручной выбор района и обычный список.

Общие фильтры --- вид услуги и диапазон цены. Дополнительно можно выбрать дату и показывать только предложения со свободным временем. Все условия применяются совместно: одна и та же услуга салона должна соответствовать виду и цене; при выборе даты для неё должен быть свободен подходящий мастер. Низкая цена другой услуги не делает салон подходящим.

Карточка показывает название, адрес, фотографию, оценку и число отзывов. При наличии точки отсчёта отображается приблизительное расстояние по прямой. Цена «от» при выбранном виде услуги означает минимум среди подходящих услуг, иначе --- минимум по каталогу с пояснением. Отсутствие отзывов обозначается текстом, а не нулевой оценкой. Список сортируется по цене подходящей услуги либо расстоянию при наличии точки отсчёта.

Карта и список сохраняют одни фильтры при переключении. Маркеры соответствуют отфильтрованным салонам выбранной области независимо от страницы списка. Перемещение карты меняет условия только по команде пользователя. При слишком большом числе маркеров система предлагает приблизить карту и явно сообщает об ограничении.

\subsection{Страницы салона и мастера}
Страница салона содержит подробное описание, адрес и точку на карте, контакты, часы работы, услуги с ценами и длительностью, список мастеров, фотографии интерьера и примеры работ, отзывы о салоне и его рейтинг. Выбор услуги открывает подходящих мастеров и доступное время.

Профиль мастера содержит имя, специализацию, описание опыта, текущий салон, выполняемые услуги, работы и отдельные отзывы о мастере. Фотография интерьера относится к салону, фотография работы --- к мастеру этого салона; работы мастеров могут показываться и на странице организации.

Цены и длительности принадлежат услугам конкретного салона. Общий вид услуги, например «Маникюр», служит для сравнения между салонами и не задаёт общую цену. В первой версии условия одной услуги салона одинаковы у всех выполняющих её мастеров.

\subsection{Запись, отмена и отзывы}
Клиент выбирает салон, услугу, мастера и дату, затем получает свободные начала с шагом 15 минут. Учитываются полная длительность, рабочие интервалы и подтверждённые записи мастера. Услуга целиком помещается в один рабочий интервал; перерывы задаются разрывами. Окончание записи может совпадать с началом следующей.

Отображённое время не удерживается за клиентом. При подтверждении сервер повторно проверяет доступность; конкурентные запросы на пересекающееся время одного мастера не должны создавать две записи. Цена и длительность фиксируются в записи и не меняются при последующем редактировании каталога.

Новая запись имеет статус «Подтверждена». До её начала клиент отменяет собственную запись, мастер --- запись к себе, владелец --- запись своего салона. Для отмены мастером или владельцем обязательна причина. Время освобождается после сохранения отмены; история остаётся. Повторная отмена не меняет результат, отмена начавшегося посещения запрещена.

После планового окончания мастер или владелец отмечает «Выполнена» либо «Неявка». Наступление времени окончания само по себе не подтверждает оказание услуги. Из «Подтверждена» допускаются переходы в «Отменена», «Выполнена» и «Неявка» при указанных условиях; конечные статусы в обычном интерфейсе не изменяются.

По выполненной собственной записи клиент независимо оставляет один отзыв о салоне и один о мастере: оценку от 1 до 5 и необязательный текст. Цель определяется через запись на сервере. Отмена и неявка не дают права на отзыв. Уникальность пары «запись--тип цели» исключает повторные отзывы.

Рейтинги салона и мастера считаются отдельно по опубликованным отзывам соответствующей цели. Оценка мастера не становится автоматически оценкой салона. Только администратор может скрыть отзыв с причиной; скрытые отзывы не входят в рейтинг и не допускают повторного создания. Редактирование отзывов и ответы на них в первую версию не входят.

\subsection{Границы учебной версии}
\begin{itemize}
\item Один город с несколькими районами и салонами, единый часовой пояс, цены в рублях.
\item Один владелец управляет одним салоном; у салона один ответственный владелец; мастер работает в одном салоне.
\item Учётная запись имеет одну роль; совмещение ролей и сети филиалов не моделируются.
\item Одна запись включает одного клиента, одного мастера и одну услугу одного салона.
\item Графики задаются на конкретные даты без повторяющихся шаблонов и переходов через полночь; горизонт записи --- 30 дней, включая текущую дату, без прошедшего времени.
\item Длительность --- положительное целое число минут и не обязана быть кратной 15; цена неотрицательна, персональные цены мастеров не вводятся.
\item Изменение графика, отключение мастера, услуги или их связи отклоняются, если нарушают будущие записи; сначала необходимо разрешить такие записи, например отменить их.
\item Перенос выполняется отменой и новой записью без гарантии сохранения прежнего времени.
\item Изображения JPEG и PNG --- до 5 МБ; до 20 активных фотографий интерьера на салон и 20 работ на мастера.
\item Онлайн-оплата, маршруты, рассылки, чаты, бонусы, групповые услуги и учёт оборудования не включаются.
\end{itemize}
Обязательный ресурсный конфликт определяется по мастеру. Параллельные записи клиента к разным мастерам отдельно не запрещаются. Допущения предназначены для ограничения учебного объёма и подлежат согласованию.

\subsection{Основные сущности и связи}
\begin{longtable}{L{3.4cm}L{11.4cm}}
\toprule Сущность & Назначение и основные сведения \\ \midrule
\endhead
Пользователь & Имя, уникальная почта, хеш пароля, роль, состояние доступа. \\
Район & Название района выбранного города; сложные географические границы не хранятся. \\
Салон & Владелец, название, описание, адрес, район, координаты, контакты, часы работы для отображения, статус публикации. \\
Вид услуги & Общий справочник сопоставимых услуг, поддерживаемый администратором. \\
Услуга салона & Салон, вид услуги, название, описание, цена, длительность, активность. \\
Мастер & Пользователь, салон, имя, специализация, опыт, активность. \\
Мастер--услуга & Уникальная пара мастера и услуги его салона. \\
Рабочее время & Мастер, дата, начало и окончание непересекающегося рабочего интервала. \\
Запись & Клиент, услуга, мастер, начало и окончание, зафиксированные название услуги, цена и длительность, статус, сведения о создании, отмене и результате. \\
Фотография & Салон, мастер для работы, тип «Интерьер» или «Работа», путь, подпись, порядок и видимость. \\
Отзыв & Запись, тип цели, оценка, текст, дата и видимость. Автор и конкретная цель определяются через запись. \\
Решение модерации & Объект, действие, администратор, время и причина скрытия, восстановления или ограничения доступа. \\
\bottomrule
\end{longtable}
Салон имеет много услуг, мастеров и фотографий. Мастер имеет много рабочих интервалов и записей. Связь мастеров и услуг --- «многие ко многим» внутри одного салона. Клиент имеет много записей, на каждую приходится не более двух отзывов разных типов.

При записи проверяется принадлежность услуги и мастера одному салону; при загрузке работы --- принадлежность мастера салону фотографии. Часы работы на странице носят справочный характер, доступность определяется рабочими интервалами, согласованными владельцем с режимом салона. Свободные варианты и рейтинги вычисляются. Объекты с историческими ссылками архивируются, а не удаляются физически.

\section{Назначение и ожидаемые результаты автоматизации}
\begin{enumerate}
\item Объединить адреса, услуги, цены, специалистов и фотографии в одном каталоге.
\item Обеспечить сравнение по месту, услуге, цене и доступности.
\item Поддержать выбор мастера через портфолио и отзывы о выполненных посещениях.
\item Автоматизировать подбор времени и предотвратить двойную запись мастера.
\item Предоставить самостоятельную отмену и историю посещений.
\item Разграничить управление платформой, салоном, расписанием мастера и данными клиента.
\end{enumerate}
Ожидается сокращение ручных согласований и повышение прозрачности выбора. Количественный экономический эффект без внедрения и измерений не заявляется.


\section{Требования к системе}
\subsection{Функциональные требования}
\begin{longtable}{L{1.6cm}L{13.2cm}}
\toprule Код & Требование \\ \midrule
\endhead
ФТ-01 & Регистрация клиента и владельца, вход и выход. Мастер получает учётную запись по приглашению владельца; администратор не регистрируется публично. \\
ФТ-02 & Владелец подаёт заявку салона с описанием, адресом и координатами. Администратор публикует её либо возвращает с причиной. \\
ФТ-03 & Публичный каталог без входа показывает зарегистрированные опубликованные активные салоны на карте и в списке с общими фильтрами. \\
ФТ-04 & Поиск поддерживает район, область карты или положение с радиусом. Отказ от геолокации не блокирует ручной поиск. \\
ФТ-05 & Фильтры вида услуги и цены применяются совместно; при выборе даты учитывается свободное время для той же услуги. Доступны страницы списка и сортировка. \\
ФТ-06 & Страница салона показывает описание, адрес, контакты, часы работы, услуги, цены, длительность, мастеров, интерьер, работы и отзывы о салоне. \\
ФТ-07 & Профиль мастера показывает специализацию, опыт, салон, услуги, работы, отзывы о мастере и переход к выбору времени. \\
ФТ-08 & Свободное время учитывает полную длительность услуги, рабочие интервалы и подтверждённые записи, исключая прошлое и пересечения. \\
ФТ-09 & Клиент подтверждает запись после просмотра её условий. Сервер повторно проверяет доступность и атомарно предотвращает конкурентные конфликты. \\
ФТ-10 & Клиент видит свои записи и отменяет будущую подтверждённую запись; время освобождается, история сохраняется. \\
ФТ-11 & Владелец ведёт только свой салон, услуги, мастеров, связи, рабочие интервалы и фотографии. Изменения, нарушающие будущие записи, отклоняются. \\
ФТ-12 & Мастер видит своё расписание и записи, редактирует собственное описание и работы; чужие данные, цены и графики ему недоступны для изменения. \\
ФТ-13 & Мастер и владелец отменяют относящиеся к ним будущие записи с причиной и отмечают выполнение либо неявку после планового окончания. \\
ФТ-14 & Клиент оставляет отдельные отзывы о салоне и мастере только по выполненной собственной записи, не более одного отзыва каждого типа. \\
ФТ-15 & Средние оценки и число опубликованных отзывов салона и мастера рассчитываются отдельно. \\
ФТ-16 & Администратор ведёт общие справочники, ограничения пользователей и публикацию салонов, модерирует отзывы и фотографии с фиксацией причин. \\
ФТ-17 & Загрузка изображений проверяет право на изменение, фактический формат, размер и количество. \\
ФТ-18 & Сообщения объясняют пустую выдачу, конфликт времени, изменение условий, недостаток прав и ошибку сохранения. Ложное подтверждение запрещено. \\
\bottomrule
\end{longtable}

\subsection{Нефункциональные требования}
Показатели являются целевыми требованиями для последующей проверки, а не результатами выполненных испытаний.
\begin{longtable}{L{1.6cm}L{13.2cm}}
\toprule Код & Требование \\ \midrule
\endhead
НФТ-01 & Изоляция салонов: сервер проверяет роль, владельца и принадлежность каждого объекта. Подмена идентификатора не открывает чужие данные. \\
НФТ-02 & Целостность: пересекающиеся записи мастера не сохраняются при конкуренции; отмена и переход статуса атомарны, уникальность отзыва контролируется БД. \\
НФТ-03 & Безопасность: адаптивные хеши паролей с солью, серверные сессии, защита от CSRF, безопасный вывод текстов без исполнения HTML; секреты не попадают в журналы. \\
НФТ-04 & Геолокация запрашивается после действия и разрешения пользователя. Текущее положение не сохраняется в профиле и истории перемещений; доступна ручная альтернатива. \\
НФТ-05 & Отказ карты не блокирует список и запись. Для геолокации требуется поддерживаемый браузером защищённый контекст. \\
НФТ-06 & Целевой 95-й процентиль серверного ответа поиска и доступности --- до 2 секунд при 20 одновременных пользователях, 100 салонах, 500 мастерах и 50\,000 записях. Внешняя карта и изображения измеряются отдельно. \\
НФТ-07 & Интерфейс на русском языке от 360 пикселей ширины, с клавиатурным доступом и текстовыми состояниями; фильтры сохраняются между представлениями. \\
НФТ-08 & Подтверждение выдаётся после фиксации транзакции. Записи, отзывы и фотографии сохраняются после перезапуска приложения. \\
НФТ-09 & Vue, Spring Boot/Spring MVC и PostgreSQL должны быть совместимы с согласованной средой демонстрации на \texttt{helios}. Версии и дисковая квота уточняются до реализации. \\
НФТ-10 & Поиск, каталог, записи, отзывы и доступ к данным разделены по ответственности внутри одного приложения; настройки окружения вынесены из кода. \\
НФТ-11 & Изображения проверяются, получают серверные имена и не исполняются; EXIF удаляется перед публикацией. Файлы хранятся отдельно от артефакта приложения. \\
\bottomrule
\end{longtable}

\section{Модели и описания основных прецедентов}
\subsection{Акторы и граница системы}
Клиент использует публичный поиск, а после входа --- собственные записи и отзывы. Мастер работает со своим профилем и расписанием. Владелец управляет своим салоном. Администратор выполняет общеплатформенные операции. Посетитель без входа --- состояние клиента до авторизации, а не отдельная бизнес-роль.

Поставщик карты является внешней технической системой. Вход служит предусловием защищённых действий. На диаграммах ассоциации обозначают участие акторов, а пунктирная стрелка «include» --- обязательную проверку при создании записи.

\begin{figure}[htbp]
\centering
\begin{tikzpicture}[font=\small,
 uc/.style={ellipse,draw,align=center,text width=5.5cm,minimum height=0.9cm}]
\node (client) at (0,0) {Клиент};
\node[uc] (a) at (5.5,2.8) {П1. Регистрация и вход};
\node[uc] (b) at (5.5,1.4) {П2. Поиск на карте и в списке};
\node[uc] (c) at (5.5,0) {П3. Просмотр салона и мастера};
\node[uc] (d) at (5.5,-1.4) {П4. Выбор времени и запись};
\node[uc] (e) at (5.5,-2.8) {П5. Свои записи и отмена};
\node[uc] (f) at (5.5,-4.2) {П6. Отзывы о посещении};
\node[uc] (g) at (5.5,-5.7) {П12. Проверка доступности};
\foreach \u in {a,b,c,d,e,f} \draw (client.east)--(\u.west);
\draw[dashed,-{Stealth}] (d.east)--++(0.7,0)|-node[pos=0.35,right]{\scriptsize include}(g.east);
\begin{scope}[on background layer]
\node[draw,fit=(a)(g),inner xsep=1.5cm,inner ysep=0.45cm,label=above:{Платформа: клиентские прецеденты}] {};
\end{scope}
\end{tikzpicture}
\caption{Модель поиска, записи и оценки посещения}
\end{figure}

\begin{figure}[htbp]
\centering
\begin{tikzpicture}[font=\small,
 uc/.style={ellipse,draw,align=center,text width=5.8cm,minimum height=1cm}]
\node[align=center] (owner) at (0,2) {Представитель\\салона};
\node (master) at (0,-1.5) {Мастер};
\node (admin) at (0,-4.5) {Администратор};
\node[uc] (a) at (6,3) {П7. Регистрация салона};
\node[uc] (b) at (6,1.4) {П8. Управление своим салоном};
\node[uc] (c) at (6,-0.2) {П9. Работа мастера};
\node[uc] (d) at (6,-1.8) {П10. Обработка посещений};
\node[uc] (e) at (6,-4.1) {П11. Администрирование\\и модерация};
\draw (owner.east)--(a.west);
\draw (owner.east)--(b.west);
\draw (owner.east)--(d.west);
\draw (master.east)--(c.west);
\draw (master.east)--(d.west);
\draw (admin.east)--(a.west);
\draw (admin.east)--(e.west);
\begin{scope}[on background layer]
\node[draw,fit=(a)(e),inner sep=0.45cm,label=above:{Платформа: профессиональные прецеденты}] {};
\end{scope}
\end{tikzpicture}
\caption{Модель работы салона, мастера и администратора}
\end{figure}
\clearpage

\subsection{П1. Регистрация и вход}
\ucfield{Акторы}{Клиент, владелец; мастер и администратор --- вход и выход.}
\ucfield{Предусловия}{Для входа учётная запись существует и не заблокирована.}
\ucfield{Триггер}{Открытие формы регистрации или входа.}
\ucfield{Основной сценарий}{Клиент или владелец вводит имя, почту и пароль. Система проверяет данные и создаёт соответствующую учётную запись. При входе открывается сессия по роли. Мастер задаёт пароль по одноразовому приглашению владельца через П8; доставка ссылки вне сайта не входит в систему.}
\ucfield{Альтернативы}{Занятая почта, неверные реквизиты, недействительное приглашение или блокировка приводят к отказу. Выход прекращает сессию.}
\ucfield{Постусловия}{Создана учётная запись или сессия; при отказе права не меняются.}

\subsection{П2. Поиск на карте и в списке}
\ucfield{Актор}{Клиент; вход не требуется.}
\ucfield{Предусловия}{Доступен публичный каталог.}
\ucfield{Триггер}{Открытие поиска или изменение фильтров.}
\ucfield{Основной сценарий}{Клиент выбирает район или область, вид услуги, цену и при необходимости дату. Система показывает подходящие карточки или маркеры. Переключение представлений сохраняет условия.}
\ucfield{Альтернативы}{По команде «Рядом со мной» после разрешения используется радиус от текущей точки. При отказе остаётся ручной поиск; при отказе карты --- список. При отсутствии совпадений предлагается изменить фильтры.}
\ucfield{Постусловия}{Показана подходящая выдача; записи не созданы.}

\subsection{П3. Просмотр салона и мастера}
\ucfield{Актор}{Клиент; вход не требуется.}
\ucfield{Предусловия}{Выбран опубликованный салон или активный мастер.}
\ucfield{Триггер}{Нажатие на карточку, маркер или имя мастера.}
\ucfield{Основной сценарий}{Система показывает описание, услуги, цены, фотографии и отзывы соответствующей цели. Клиент выбирает услугу и мастера, переходит к подбору времени.}
\ucfield{Альтернативы}{Отсутствие фотографий или отзывов явно обозначается. Для снятого с публикации салона новые записи недоступны, но существующая запись остаётся доступной её участникам.}
\ucfield{Постусловия}{Клиент получил сведения для выбора.}

\subsection{П4. Выбор времени и создание записи}
\ucfield{Актор}{Клиент.}
\ucfield{Предусловия}{Выбраны опубликованный салон, активная услуга и подходящий мастер; подтверждение требует входа.}
\ucfield{Триггер}{Выбор даты посещения.}
\ucfield{Основной сценарий}{Система показывает свободное время. Клиент выбирает начало, проверяет салон, услугу, мастера, окончание и цену. После подтверждения сервер выполняет П12 в транзакции, сохраняет запись и возвращает её номер.}
\ucfield{Альтернативы}{Занятое время отклоняется, варианты обновляются. Изменение цены или длительности требует нового подтверждения. При потере ответа клиент проверяет свои записи; повторный запрос не создаёт пересекающийся дубликат.}
\ucfield{Постусловия}{Сохранена одна допустимая запись либо новая запись отсутствует.}

\subsection{П5. Свои записи и отмена}
\ucfield{Актор}{Клиент.}
\ucfield{Предусловия}{Клиент авторизован.}
\ucfield{Триггер}{Открытие раздела «Мои записи».}
\ucfield{Основной сценарий}{Система показывает будущие и прошлые собственные записи с условиями и статусом. При необходимости клиент выбирает будущую запись и подтверждает отмену. Сервер проверяет владельца и состояние и сохраняет отмену.}
\ucfield{Альтернативы}{Отмена чужой или начавшейся записи запрещена. Повторная отмена возвращает текущее состояние. Просмотр не требует отмены.}
\ucfield{Постусловия}{При отмене время освобождено, история сохранена.}

\subsection{П6. Отзывы о посещении}
\ucfield{Актор}{Клиент.}
\ucfield{Предусловия}{Собственная запись имеет статус «Выполнена».}
\ucfield{Триггер}{Выбор оценки салона или мастера в истории.}
\ucfield{Основной сценарий}{Клиент выбирает тип цели, оценку и текст. Сервер определяет цель по записи, проверяет уникальность и публикует отзыв, обновляя рейтинг соответствующей цели.}
\ucfield{Альтернативы}{Повторный тип отзыва, чужая запись, отмена или неявка приводят к отказу. Другой тип отзыва по той же выполненной записи допускается.}
\ucfield{Постусловия}{Сохранён отзыв, связанный с выполненным посещением.}


\subsection{П7. Регистрация и публикация салона}
\ucfield{Акторы}{Владелец, администратор.}
\ucfield{Предусловия}{Владелец авторизован и не управляет другим салоном.}
\ucfield{Триггер}{Создание страницы организации.}
\ucfield{Основной сценарий}{Владелец вводит описание, адрес, район, контакты и часы работы, отмечает координаты. Система сохраняет заявку. Администратор проверяет и публикует её.}
\ucfield{Альтернативы}{Неполные данные отклоняются. Возвращённую заявку можно исправить и подать повторно. Салон без услуг или графика виден после публикации, но не имеет времени для записи.}
\ucfield{Постусловия}{Салон опубликован либо остаётся непубличной заявкой.}

\subsection{П8. Управление своим салоном}
\ucfield{Актор}{Владелец.}
\ucfield{Предусловия}{Выполнен вход; салон принадлежит владельцу.}
\ucfield{Триггер}{Открытие кабинета салона.}
\ucfield{Основной сценарий}{Владелец редактирует страницу, услуги, цены и фотографии, приглашает мастеров и назначает услуги и рабочие интервалы. Сервер проверяет принадлежность, данные и влияние на будущие записи, затем сохраняет изменения.}
\ucfield{Альтернативы}{Чужая услуга, пересечение графиков, недопустимый файл и нарушение будущих записей отклоняются. Приглашение на уже занятую почту не меняет роль существующей учётной записи и отклоняется.}
\ucfield{Постусловия}{Собственный каталог обновлён; чужие данные и согласованные условия записей сохранены.}

\subsection{П9. Работа мастера}
\ucfield{Актор}{Мастер.}
\ucfield{Предусловия}{Выполнен вход, мастер связан с салоном.}
\ucfield{Триггер}{Открытие личного кабинета.}
\ucfield{Основной сценарий}{Мастер выбирает дату и видит рабочие интервалы, свои записи, имена клиентов, услуги и время. При необходимости редактирует своё описание и фотографии работ.}
\ucfield{Альтернативы}{Пустое расписание обозначается сообщением. Изменение чужого профиля, цены или графика запрещено; недопустимый файл отклоняется.}
\ucfield{Постусловия}{Показано собственное расписание, разрешённые изменения сохранены.}

\subsection{П10. Обработка посещений}
\ucfield{Акторы}{Мастер, владелец.}
\ucfield{Предусловия}{Запись относится к текущему мастеру или салону владельца.}
\ucfield{Триггер}{Выбор действия в журнале.}
\ucfield{Основной сценарий}{Мастер видит свои записи, владелец --- журнал своего салона с фильтрами по дате, мастеру, клиенту и статусу. После планового окончания актор отмечает выполнение. Сервер проверяет права и статус, сохраняет результат и инициатора. Клиенту становится доступна оценка.}
\ucfield{Альтернативы}{После окончания можно отметить неявку без права на отзыв. До начала можно отменить с причиной. Неподходящее время или конечный статус запрещают переход.}
\ucfield{Постусловия}{Состояние и история согласованы; право на отзыв соответствует результату.}

\subsection{П11. Администрирование и модерация}
\ucfield{Актор}{Администратор платформы.}
\ucfield{Предусловия}{Выполнен административный вход.}
\ucfield{Триггер}{Выбор справочника или объекта модерации.}
\ucfield{Основной сценарий}{Администратор ведёт районы и виды услуг, проверяет контент. При нарушении скрывает отзыв или фотографию, ограничивает пользователя или публикацию салона с сохранением причины и времени.}
\ucfield{Альтернативы}{Снятие салона с публикации запрещает новые записи, но не отменяет старые: владелец сохраняет доступ к их урегулированию. При блокировке владельца администратор может в исключительном порядке отменить затронутые будущие записи с причиной. Блокировка мастера запрещает новые записи к нему, а существующие обрабатывает владелец. Восстановление доступа и контента фиксируется отдельным решением.}
\ucfield{Постусловия}{Видимость и права обновлены, скрытые отзывы исключены из рейтинга, история сохранена.}

\subsection{П12. Проверка доступности перед сохранением}
\ucfield{Контекст}{Обязательная включаемая часть П4.}
\ucfield{Предусловия}{Открыта транзакция и обеспечена синхронизация изменений выбранного мастера.}
\ucfield{Основной сценарий}{Сервер проверяет публикацию салона, активность и принадлежность услуги и мастера, их связь, актуальность условий, дату, шаг начала, полное вхождение в рабочий интервал и отсутствие пересечений.}
\ucfield{Альтернативы}{Любое нарушение останавливает сохранение с конкретной причиной.}
\ucfield{Постусловия}{Разрешение действует только в текущей транзакции и не заменяется предварительной проверкой интерфейса.}

\subsection{Связь требований и прецедентов}
\begin{center}
\begin{tabular}{ll}
\toprule Требования & Прецеденты \\ \midrule
ФТ-01 & П1, П8 \\
ФТ-02 & П7 \\
ФТ-03--ФТ-05 & П2 \\
ФТ-06, ФТ-07 & П3 \\
ФТ-08, ФТ-09 & П4, П12 \\
ФТ-10 & П5 \\
ФТ-11 & П8 \\
ФТ-12 & П9 \\
ФТ-13 & П10 \\
ФТ-14, ФТ-15 & П6, П11 \\
ФТ-16 & П7, П11 \\
ФТ-17 & П8, П9 \\
ФТ-18 & П1--П12 \\
\bottomrule
\end{tabular}
\end{center}

\section{Предлагаемая архитектура}
\subsection{Общий подход и технологии}
Предлагается клиент-серверная архитектура с монолитной серверной частью. Несколько салонов обслуживаются одним приложением и общей PostgreSQL, а принадлежность данных проверяется сервером. Для учебного масштаба не требуются отдельные микросервисы и поисковый кластер.
\begin{itemize}
\item \textbf{Vue 3} реализует каталог, фильтры, страницы салонов и мастеров и кабинеты по ролям.
\item \textbf{Leaflet} отображает карту, маркеры и выбранную область; подложка подключается от согласованного поставщика \cite{leaflet}. Салоны поступают из собственной БД, а не из внешнего каталога карты.
\item \textbf{Spring Boot и Spring MVC} обеспечивают Java-приложение и HTTP API; Spring Security --- сессии и права, Spring Data JPA --- доступ к данным.
\item \textbf{PostgreSQL} хранит пользователей, салоны, каталог, графики, записи, отзывы и метаданные изображений. Координаты хранятся числовыми полями; для учебной выборки расстояния вычисляются без обязательного PostGIS.
\item \textbf{Файловое хранилище} в отдельном постоянном каталоге содержит изображения; их пути и связи хранятся в БД. Обновление JAR не удаляет фотографии.
\end{itemize}
Геолокация браузера требует разрешения пользователя и защищённого контекста \cite{geolocation}. Поставщик подложки согласуется отдельно. При использовании публичных тайлов OpenStreetMap соблюдаются требования атрибуции и правила сервиса; не предполагается неограниченная гарантированная доступность \cite{osm}. Адрес вводится владельцем, координаты задаются точкой на карте, поэтому внешний геокодер не является обязательным компонентом.

Серверная часть должна демонстрироваться на \texttt{helios}. Статические файлы Vue могут обслуживаться тем же Spring Boot-приложением. Версии Java и Spring Boot, защищённый браузерный доступ и место для изображений уточняются до реализации. Факт успешного развёртывания не заявляется.

\begin{figure}[htbp]
\centering
\begin{tikzpicture}[font=\small,
 box/.style={draw,rounded corners,align=center,text width=10.8cm,minimum height=1.1cm},
 arrow/.style={-{Stealth},thick}]
\node[box] (ui) {Браузер: Vue, Leaflet, разрешённая геолокация\\Клиент, мастер, владелец, администратор};
\node[box,below=0.8cm of ui] (api) {Spring Boot / Spring MVC на \texttt{helios}\\HTTP API и Spring Security};
\node[box,below=0.8cm of api] (logic) {Каталог и поиск, запись, кабинеты, отзывы, фотографии\\Проверка роли и принадлежности данных салону};
\node[box,below=0.8cm of logic] (data) {PostgreSQL: постоянные данные и транзакции\\Отдельный каталог файлов: фотографии};
\draw[arrow] (ui)--node[right]{JSON} (api);
\draw[arrow] (api)--(logic);
\draw[arrow] (logic)--(data);
\end{tikzpicture}
\caption{Основные уровни платформы; подложка карты загружается браузером из внешнего сервиса}
\end{figure}

\subsection{Ответственность компонентов}
Поиск обрабатывает географию, услугу, цену и дату как единый запрос. Один серверный механизм формирует список и маркеры, чтобы их результаты не расходились. При небольшом объёме достаточно фильтрации в PostgreSQL и вычисления расстояний по координатам; внешняя карта отвечает за отображение.

Сервис записи принимает решение о доступности и выполняет проверку и сохранение в транзакции с блокировкой строки мастера. Изменение графика и отмена используют ту же синхронизацию. Изменения услуги и публикации салона также согласуются с созданием записи; для блокировок устанавливается единый порядок. Детальный протокол определяется на этапе реализации. Проверка только в браузере недостаточна.

Сервис отзывов проверяет владельца и выполнение записи, обеспечивает уникальность типа отзыва и рассчитывает отдельные рейтинги. Сервис фотографий проверяет загрузку и формирует безопасные имена; неудачное сохранение не должно оставлять публичную ссылку на отсутствующий файл.

Сервис доступа различает личные данные клиента, данные мастера и данные салона. Идентификатор салона из запроса не доказывает полномочия. Профессиональные кабинеты получают только необходимые сведения о связанных посещениях, без закрытых реквизитов учётных записей.

Группы HTTP API: учётные записи и сессии; поиск; страницы салонов и мастеров; услуги и доступность; записи и статусы; отзывы; фотографии; профессиональные кабинеты; модерация. Конкретные адреса и структуры запросов определяются на следующем этапе.

\section{Заключение}
Новая постановка описывает платформу поиска и сравнения нескольких салонов с онлайн-записью. Определены четыре роли, единая выдача на карте и в списке, фильтры, страницы организаций и мастеров, фотографии, отзывы по выполненным посещениям и отмена записей.

Модель связывает цены с услугами конкретного салона, разделяет отзывы о салоне и мастере и сохраняет проверку конфликтов времени как обязательное условие записи. Архитектура использует Vue, Spring Boot/Spring MVC и PostgreSQL. Следующий шаг --- согласование объёма, прецедентов и среды демонстрации с преподавателем.

\begin{thebibliography}{9}
\addcontentsline{toc}{section}{Список источников}
\bibitem{leaflet} Leaflet. Официальная документация библиотеки интерактивных карт. URL: \url{https://leafletjs.com/} (дата обращения: 15.09.2026).
\bibitem{geolocation} MDN Web Docs. Браузерный интерфейс геолокации. URL: \url{https://developer.mozilla.org/en-US/docs/Web/API/Geolocation_API} (дата обращения: 15.09.2026).
\bibitem{osm} OpenStreetMap Foundation. Правила использования картографических тайлов. URL: \url{https://operations.osmfoundation.org/policies/tiles/} (дата обращения: 15.09.2026).
\end{thebibliography}

